% ============================================================
%  SY02-19 《几何朗兰兹的一个算术应用》译稿导言区
%  底本：Sam Raskin, An Arithmetic Application of Geometric Langlands,
%        Proc. Sympos. Pure Math. 112.1 (2025), 42 pp.
%  体例：中文数学论文；定理类编号「节.小节.子小节.序号」，公式编号「节.序号」
% ============================================================

\usepackage{geometry}
\geometry{a4paper, left=2.7cm, right=2.7cm, top=2.8cm, bottom=2.7cm}
\linespread{1.22}
\setlength{\parskip}{0.25em}
\setlength{\emergencystretch}{2em}
\sloppy
\hyphenpenalty=10000
\tolerance=3000

% ---------- 数学 ----------
\usepackage{amsmath}
\usepackage{amsthm}
\usepackage{amssymb}
\usepackage{amsfonts}
\usepackage{mathtools}
\usepackage{mathrsfs}
\usepackage{stmaryrd}
\usepackage{bm}
\allowdisplaybreaks[2]

% ---------- 图表 ----------
\usepackage{graphicx}
\usepackage{xcolor}
\usepackage{tikz}
\usetikzlibrary{cd, arrows.meta, positioning, calc}
\usepackage[most]{tcolorbox}
\usepackage{enumitem}
\setlist{itemsep=2pt, topsep=3pt, parsep=0pt}

% ---------- 超链接 ----------
\usepackage[hyperfootnotes=false, colorlinks=true, linkcolor=blue!45!black, citecolor=green!35!black,
            urlcolor=teal!60!black, bookmarksnumbered=true]{hyperref}
\usepackage{bookmark}

% ---------- 中文微调 ----------
\xeCJKsetup{CJKecglue={\hskip 0.24em plus 0.05em minus 0.05em}}
\ctexset{
  section    = {name={,}, number=\arabic{section},
                aftername={.\enspace}, format=\Large\bfseries\centering},
  subsection = {number=\arabic{section}.\arabic{subsection},
                aftername={.\enspace}, format=\large\bfseries},
  subsubsection = {number=\arabic{section}.\arabic{subsection}.\arabic{subsubsection},
                aftername={.\enspace}, format=\normalsize\bfseries}
}

% ---------- 定理环境：共用计数器，按子小节重置（复现 1.2.4.1、4.3.3.1 体例） ----------
\theoremstyle{plain}
\newtheorem{thm}{定理}[subsubsection]
\newtheorem{prop}[thm]{命题}
\newtheorem{lem}[thm]{引理}
\newtheorem{cor}[thm]{推论}
\newtheorem{conj}{猜想}[subsection]
\newtheorem{notation}[thm]{记号}
\newtheorem{question}[thm]{问题}
\newtheorem{hypothesis}[thm]{假设}
\theoremstyle{definition}
\newtheorem{defn}[thm]{定义}
\newtheorem{example}[thm]{例}
\newtheorem{convention}[thm]{约定}
\theoremstyle{remark}
\newtheorem{remark}[thm]{注}
\newtheorem{warning}[thm]{警告}

% 定理 A 是原文专名，不参与自动编号
\newtheorem*{thmA}{定理 A}

% 原文把「1.1.1. 背景」这类编号段落作为正文编号，而非标题；
% \pnum 递增 subsubsection 计数器并run-in 输出，定理编号随之对齐
\newcommand{\pnum}[1]{\par\refstepcounter{subsubsection}%
  \noindent\textbf{\thesubsubsection.}\ #1\ }
\newcommand{\pnumt}[2]{\par\refstepcounter{subsubsection}%
  \noindent\textbf{\thesubsubsection.}\ \textbf{#1}\ #2\ }

\numberwithin{equation}{section}
\renewcommand{\proofname}{证明}
\renewcommand{\qedsymbol}{$\square$}
\renewcommand{\refname}{参考文献}
\renewcommand{\contentsname}{目\quad 录}

% ---------- 译注 ----------
\usepackage{manyfoot}
\DeclareNewFootnote[plain]{trans}
\renewcommand{\thefootnotetrans}{译\arabic{footnotetrans}}
\newcommand{\zy}[1]{\footnotetrans{#1}}
\newtcolorbox{yizhu}[1][]{
  enhanced, breakable, colback=orange!4!white, colframe=orange!55!black,
  boxrule=0.5pt, left=6pt, right=6pt, top=4pt, bottom=4pt,
  fonttitle=\bfseries\small, title=#1,
  before skip=6pt, after skip=6pt
}

% ---------- 本文记号 ----------
% ============================================================
%  SY02-19 记号宏（按原文用法整理）
% ============================================================

% 常用缩写
\newcommand{\LS}{\mathrm{LS}}
\newcommand{\Bun}{\mathrm{Bun}}
\newcommand{\Shv}{\mathrm{Shv}}
\newcommand{\IndCoh}{\mathrm{IndCoh}}
\newcommand{\QCoh}{\mathrm{QCoh}}
\newcommand{\Coh}{\mathrm{Coh}}
\newcommand{\Perf}{\mathrm{Perf}}
\newcommand{\Rep}{\mathrm{Rep}}
\newcommand{\Nilp}{\mathrm{Nilp}}
\newcommand{\qLisse}{\mathrm{qLisse}}
\newcommand{\Eis}{\mathrm{Eis}}
\newcommand{\Frob}{\mathrm{Frob}}
\newcommand{\Dmod}{\mathrm{D}\text{-}\mathrm{mod}}
\newcommand{\DGCat}{\mathrm{DGCat}}
\newcommand{\Vect}{\mathrm{Vect}}
\newcommand{\Spectra}{\mathrm{Spectra}}
\newcommand{\Exc}{\mathrm{Exc}}
\newcommand{\Sht}{\mathrm{Sht}}
\newcommand{\Eisspec}{\mathrm{Eis}^{\mathrm{spec}}}
\newcommand{\Eissoyt}{\mathrm{Eis}^{\mathrm{toy}}}
\newcommand{\arthm}{\mathrm{arthm}}
\newcommand{\restr}{\mathrm{restr}}
\newcommand{\cusp}{\mathrm{cusp}}
\newcommand{\unr}{\mathrm{unr}}
\newcommand{\nonres}{\mathrm{non}\text{-}\mathrm{res}}
\newcommand{\triv}{\mathrm{triv}}
\newcommand{\op}{\mathrm{op}}
\newcommand{\et}{\text{ét}}
\newcommand{\dR}{\mathrm{dR}}

% 算子
\DeclareMathOperator{\Aut}{Aut}
\DeclareMathOperator{\Hom}{Hom}
\DeclareMathOperator{\Ext}{Ext}
\DeclareMathOperator{\Spec}{Spec}
\DeclareMathOperator{\ad}{ad}
\DeclareMathOperator{\Ad}{Ad}
\DeclareMathOperator{\ch}{ch}
\DeclareMathOperator{\vol}{vol}
\DeclareMathOperator{\cl}{cl}
\DeclareMathOperator{\tr}{tr}
\DeclareMathOperator{\id}{id}
\DeclareMathOperator{\Sect}{Sect}
\DeclareMathOperator{\colim}{colim}

% 常用对象
\newcommand{\bG}{\check{G}}
\newcommand{\bT}{\check{T}}
\newcommand{\bB}{\check{B}}
\newcommand{\bN}{\check{N}}
\newcommand{\bg}{\check{\mathfrak g}}
\newcommand{\bt}{\check{\mathfrak t}}
\newcommand{\hG}{\widehat{G}}
\newcommand{\hT}{\widehat{T}}
\newcommand{\hB}{\widehat{B}}
\newcommand{\hZ}{\widehat{Z}}
\newcommand{\hf}{\widehat{f}}
\newcommand{\fh}{\widehat{h}}
\newcommand{\wt}{\widetilde}
\newcommand{\Om}{\Omega}
\newcommand{\om}{\omega}
\newcommand{\e}{\mathrm{e}}
\newcommand{\F}{\mathbb{F}}
\newcommand{\Z}{\mathbb{Z}}
\newcommand{\Q}{\mathbb{Q}}
\newcommand{\A}{\mathbb{A}}
\newcommand{\OO}{\mathcal{O}}
\newcommand{\ZZ}{\mathbb{Z}}
\newcommand{\QQ}{\mathbb{Q}}
\newcommand{\Fq}{\F_q}
\newcommand{\Fb}{\overline{\F}_q}

